\section{Appendix}

\subsection{Source Code Listings}

\subsubsection{Backend API Logic (route.ts)}
The following TypeScript code demonstrates the core logic for the \texttt{/api/analyze} endpoint, including the prompt construction and OpenAI API integration.

\begin{lstlisting}[language=JavaScript, caption=Next.js API Route for Stock Analysis]
import { NextResponse } from 'next/server';
import OpenAI from 'openai';

const openai = new OpenAI({
    apiKey: process.env.OPENAI_API_KEY,
});

export async function POST(request: Request) {
    try {
        const { symbol, historical, quote } = await request.json();

        if (!symbol || !historical || !quote) {
            return NextResponse.json(
                { error: 'Missing required data' }, 
                { status: 400 }
            );
        }

        // Prepare data for prompt (limit to last 30 days)
        const recentHistory = historical.slice(-30).map((day: any) => ({
            date: day.date,
            close: day.close,
            volume: day.volume,
        }));

        const prompt = `
      Analyze the following stock data for ${symbol}.
      Current Price: ${quote.regularMarketPrice}
      Recent History (last 30 days): ${JSON.stringify(recentHistory)}

      Based on this data, determine if the trend suggests a 
      "BUY" or "SELL" or "HOLD".
      Provide a recommendation and a short textual explanation 
      (max 3 sentences) explaining the trend.
      
      Format the response as JSON:
      {
        "recommendation": "BUY" | "SELL" | "HOLD",
        "explanation": "string"
      }
    `;

        const completion = await openai.chat.completions.create({
            messages: [{ role: 'user', content: prompt }],
            model: 'gpt-4o',
            response_format: { type: "json_object" },
        });

        const content = completion.choices[0].message.content;
        if (!content) throw new Error('No content from OpenAI');
        
        const result = JSON.parse(content);
        return NextResponse.json(result);
    } catch (error) {
        console.error('Error analyzing stock:', error);
        return NextResponse.json(
            { error: 'Failed to analyze stock' }, 
            { status: 500 }
        );
    }
}
\end{lstlisting}

\subsubsection{Frontend Component (StockChart.tsx)}
This component renders the visualization of the stock data, crucial for the user's verification of the AI's analysis.

\begin{lstlisting}[language=JavaScript, caption=React Component for Stock Charting]
'use client';

import { 
  LineChart, Line, XAxis, YAxis, CartesianGrid, 
  Tooltip, ResponsiveContainer 
} from 'recharts';

export default function StockChart({ data }: { data: any[] }) {
  return (
    <div className="h-64 w-full bg-gray-800/50 rounded-xl p-4">
      <ResponsiveContainer width="100%" height="100%">
        <LineChart data={data}>
          <CartesianGrid strokeDasharray="3 3" stroke="#374151" />
          <XAxis 
            dataKey="date" 
            stroke="#9CA3AF" 
            tick={{fontSize: 12}}
            tickFormatter={(str) => new Date(str).toLocaleDateString()}
          />
          <YAxis 
            stroke="#9CA3AF" 
            domain={['auto', 'auto']}
            tick={{fontSize: 12}}
          />
          <Tooltip 
            contentStyle={{
                backgroundColor: '#1F2937', 
                border: 'none', 
                borderRadius: '8px'
            }}
          />
          <Line 
            type="monotone" 
            dataKey="close" 
            stroke="#3B82F6" 
            strokeWidth={2} 
            dot={false} 
          />
        </LineChart>
      </ResponsiveContainer>
    </div>
  );
}
\end{lstlisting}

\subsection{Prompt Engineering Iterations}
The prompt used in the final system was the result of several iterations. Here we document the evolution of the prompt to demonstrate the importance of specific instructions.

\subsubsection{Iteration 1: Naive Prompt}
\begin{quote}
    "Should I buy this stock? Here is the data: [JSON]"
\end{quote}
\textbf{Result}: The model often refused to answer, citing "I am not a financial advisor." When it did answer, the reasoning was vague.

\subsubsection{Iteration 2: Role-Playing}
\begin{quote}
    "Act as a financial advisor. Analyze this data [JSON] and tell me to buy or sell."
\end{quote}
\textbf{Result}: Better, but the model was too aggressive, often recommending "BUY" on weak signals. It also failed to return valid JSON.

\subsubsection{Iteration 3: Final Structured Prompt}
\begin{quote}
    "Analyze the following stock data... Format the response as JSON... Provide a recommendation... (max 3 sentences)..."
\end{quote}
\textbf{Result}: Consistent JSON output, conservative "HOLD" bias, and clear explanations.
